% !TeX TXS-program:compile = txs:///arara
% arara: lualatex: {shell: yes, synctex: no, interaction: batchmode}
% arara: lualatex: {shell: yes, synctex: no, interaction: batchmode} if found('log', '(undefined references|Please rerun|Rerun to get)')

\documentclass[a4paper]{scrartcl}
\usepackage{mathtools, array}
\usepackage[british]{babel}
\usepackage{libertinus-otf}
\usepackage[Scale=MatchLowercase]{fantasquesansmono-otf}
\usepackage{unicodefonttable}
\usepackage{framed}
\usepackage{listings}
\usepackage{minted-code}
\usepackage{xcolor}
\usepackage{realscripts}
\usepackage{graphicx}
\usepackage{microtype}

\newcommand*{\pkg}[1]{\texttt{#1}}
\newcommand*{\file}[1]{\texttt{#1}}
\newcommand*{\ttbold}[1]{\texttt{\textbf{#1}}}
\newcommand*{\opt}[1]{\texttt{#1}}
\newcommand*{\cmd}[1]{\texttt{\textbackslash #1}}\newcommand*{\showtchar}[1]{\cmd{#1}~\csname #1\endcsname}
\newcommand*{\showmchar}[1]{\cmd{#1}~$(\csname #1\endcsname)$}
\newcommand*{\showmchardollar}[1]{\texttt{\$\cmd{#1}\$}~$(\csname #1\endcsname)$}

\newcommand{\abc}{abcdefghijklmnopqrstuvwxyz}
\newcommand{\ABC}{ABCDEFGHIJKLMNOPQRSTUVWXYZ}
\newcommand{\samplenumbers}{0123456789}

\newcommand\samplettxt{oO08 iIlL1 g9qCGQ <=> -> != !== \string~\string~>}
\newcommand\samplett[1][]{\cmd{#1} & \fbox{\csname#1\endcsname\samplettxt} \\}

\setlength{\parindent}{0pt}
\setlength{\parskip}{6pt plus 2pt minus 2pt}

\usepackage{hyperref}
\hypersetup{pdftitle={FantasqueSansMono User’s Guide for LaTeX},
	pdfauthor={Cédric PIERQUET},
	bookmarksopen,
	colorlinks
}
\newcommand*{\hlabel}[1]{\phantomsection\label{#1}}

\title{FantasqueSansMono family fonts (\textit{experimental}) \\User’s Guide for LaTeX (v1.8.0)}
\author{Cédric Pierquet \\ \texttt{cpierquet@outlook.fr}}

\date{\today}
\newcommand*{\version}{1.8.0}

\begin{document}

\maketitle

\section{What is FantasqueSansMono ?}

\subsection{Presentation (actual version is v1.8.0)}

A programming font, designed with functionality in mind, and with some wibbly-wobbly handwriting-like fuzziness that makes it unassumingly cool.

\medskip

For further information, \url{https://github.com/belluzj/fantasque-sans} is the best source!

\subsection{License}

Fonts are licensed under the SIL Open Font License, Version 1.1.

They require LuaTeX or XeTeX as engine and the \pkg{fontspec} package\footnote{Please read the documentation \file{fontspec.pdf}.}.

\pagebreak

\section{Usage}

\subsection{Loading mono/code font}

Files \file{FantasqueSansCode.fontspec} and \file{FantasqueSansMono.fontspec} are provided to ensure that \texttt{Regular}, \texttt{\textit{Italic}}, \texttt{\textbf{Bold}}, \texttt{\textit{\textbf{BoldItalic}}} are properly loaded.

\smallskip

A basic call for FantasqueSansMono mono font could be:

\begin{codeblock}[minted language=latex, title=Loading]
\usepackage{fantasquesansmono-otf}
\end{codeblock}

It loads \opt{FantasqueSansMono} as mono font.

\medskip

%\begin{lstlisting}[language=python,basicstyle=\footnotesize\ttfamily,commentstyle=\itshape\color{gray},keywordstyle=\bfseries\color{magenta},tabsize=4,frame=single]
\begin{codeblock}[minted language=python3, title=Sample python code]
def Fibonacci(n) :
  # Check if input is 0 then it will print incorrect input
  if n < 0 :
    print("Incorrect input")
  elif n == 0 :
    return 0
  elif n == 1 or n == 2 :
    return 1
  else :
    return Fibonacci(n-1) + Fibonacci(n-2)
\end{codeblock}
%\end{lstlisting}

\subsection{Options}

Options can be given within the package:

\begin{itemize}
	\item the boolean \opt{nott} for not loading fantasquesansmono as mono font;
  \item the boolean \opt{code} for loading ligatures;
	\item \opt{StylisticSet=...}
	\item \opt{RawFeature=...}
	\item \opt{Scale=...}
\end{itemize}

\subsection{Manual loading}

It's also possible to load mono font with \cmd{setmonofont}, thanks to \opt{fontspec} files for example.

\begin{codeblock}[minted language=latex, title=Manual loading]
%w/o package loaded
\usepackage{fontspec}
\setmonofont{FantasqueSansMono}[options]
\setmonofont{FantasqueSansCode}[options]
\end{codeblock}

\subsection{Fontfamily}

Mono/Code \opt{fontfamily} are provided:

\begin{tabular}{ll}
	\samplett[fantasquesansmono]
	\samplett[fantasquesansmonocode]
\end{tabular}

\opt{code} version is declared with ligatures.

\smallskip

Only \opt{StylisticSet={1}} is avalaible, with alternate \opt{k}/\opt{\addfontfeatures{StylisticSet={1}}k} version.

\section{Samples}

\subsection{Mono}

%\begin{lstlisting}[language=TeX,basicstyle=\footnotesize\fantasquesansmono,commentstyle=\itshape\color{gray},keywordstyle=\bfseries\color{blue},tabsize=4,frame=single,columns=flexible,showstringspaces=false]
%\lstset{
%  language=python,
%  basicstyle=\footnotesize\fantasquesansmono,
%  commentstyle=\itshape\color{gray},
%  keywordstyle=\bfseries\color{magenta},
%  tabsize=4,
%  frame=single,
%  columns=flexible,
%  showstringspaces=false
%}
%\end{lstlisting}

\begin{codeblock}[minted language=java, title=Sample java code]
// a sample of java code

const similar = "oO08 iIlL1 g9qCGQ <=> -> != !== ~~>"
const diacritics_etc = "â é ù ï ø ç à Ē Æ œ"

window.toggleFavorite = (alias) => {
  try {
    let favorites = JSON.parse(localStorage.getItem('favorites')) || []
    if (favorites.indexOf(alias) > -1) {
      favorites = favorites.filter((v) => {
        return v !== alias
      })
    } else {
      favorites.push(alias)
    }
    localStorage.setItem('favorites', JSON.stringify(Array.from(new Set(favorites))))
  } catch (err) {
    // eslint-disable-next-line no-console
    console.error('could not save favorite', err)
  }
  renderSelectList()
  return false
}
\end{codeblock}

\begin{codeblock}[minted language=python3, title=Sample python code]
#a sample of python code
def Fibonacci(n: int) -> int :
  # Check if input is 0 then it will print incorrect input
  if n < 0 :
    print("Incorrect input")
  elif n == 0 :
    return 0
  elif 1 <= n <= 2 :
    return 1
  else :
    return Fibonacci(n-1) + Fibonacci(n-2)
\end{codeblock}

\subsection{Code}

\begin{codeblock}[minted language=java, title=Sample java code with ligatures,minted options app={formatcom=\fantasquesansmonocode}]
// a sample of java code

const similar = "oO08 iIlL1 g9qCGQ <=> -> != !== ~~>"
const diacritics_etc = "â é ù ï ø ç à Ē Æ œ"

window.toggleFavorite = (alias) => {
  try {
    let favorites = JSON.parse(localStorage.getItem('favorites')) || []
    if (favorites.indexOf(alias) > -1) {
      favorites = favorites.filter((v) => {
        return v !== alias
      })
    } else {
      favorites.push(alias)
    }
    localStorage.setItem('favorites', JSON.stringify(Array.from(new Set(favorites))))
  } catch (err) {
    // eslint-disable-next-line no-console
    console.error('could not save favorite', err)
  }
  renderSelectList()
  return false
}
\end{codeblock}

\begin{codeblock}[minted language=python3, title=Sample python code with ligatures,minted options app={formatcom=\fantasquesansmonocode}]
#a sample of python code
def Fibonacci(n: int) -> int :
  # Check if input is 0 then it will print incorrect input
  if n < 0 :
    print("Incorrect input")
  elif n == 0 :
    return 0
  elif 1 <= n <= 2 :
    return 1
  else :
    return Fibonacci(n-1) + Fibonacci(n-2)
\end{codeblock}

\pagebreak

\section{List of glyphs}

\displayfonttable{FantasqueSansMono-Regular}

\end{document}